As discussed previously, the baseline model’s performance is significantly constrained by two major limitations in the dataset 
and problem structure: severe class imbalance and the ordinal nature of diabetic retinopathy severity grades. The dataset is 
heavily skewed toward lower-grade samples, causing the model to bias predictions toward majority classes while underperforming 
on minority classes representing more severe disease stages. In addition, diabetic retinopathy classification is inherently ordinal,
meaning that misclassifications should not all be treated equally. For example, predicting a Grade 4 image as Grade 3 is considerably
less severe than predicting it as Grade 0, yet standard Cross Entropy Loss penalizes both errors similarly because it treats
classes as independent categories rather than ordered progression levels. This limits the model’s ability to learn the structured 
severity relationship between grades.

To address these shortcomings, we propose the use of a hybrid loss function that combines Focal Loss and CORAL 
(Consistent Rank Logits) Loss. Focal Loss is specifically designed to mitigate class imbalance by dynamically down-weighting 
easy, high-confidence predictions and assigning greater emphasis to hard or misclassified examples. This will ensure that
minority and difficult samples contribute more significantly during backpropagation, improving sensitivity to
underrepresented disease stages. CORAL Loss, on the other hand, is designed for ordinal classification problems
by transforming multi-class prediction into a series of ordered binary threshold tasks. Instead of treating grades
as unrelated labels, CORAL will learn rank-consistent decision boundaries that explicitly preserve severity progression,
allowing the model to better understand the ordered relationship between diabetic retinopathy grades. By integrating
Focal weighting into the CORAL framework, the proposed loss function will simultaneously address both imbalance and ordinality,
enabling the model to prioritize difficult minority samples while preserving clinically meaningful grading structure. The 
following is the loss function we formulated:

\[
\mathcal{L}_{\text{Focal-CORAL}} =
- \sum_{i=1}^{N}\sum_{k=1}^{K-1}
\alpha_k (1-p_{t,i}^{(k)})^\gamma \log\left(p_{t,i}^{(k)}\right)
\]

\[
p_{t,i}^{(k)} =
\begin{cases}
p_i^{(k)}, & \text{if } y_i^{(k)} = 1 \\
1 - p_i^{(k)}, & \text{if } y_i^{(k)} = 0
\end{cases}
\]

\[
p_i^{(k)} = \sigma(z_i^{(k)})
\]

\[
y_i^{(k)} = \mathbb{I}(y_i > k), \quad k = 1,2,\dots,K-1
\]

In addition to loss function improvements, extensive preprocessing will be applied to enhance image quality and reduce
irrelevant variability in retinal fundus images. Medical retinal images often suffer from inconsistent illumination,
poor contrast, noise, and imaging artifacts due to differences in acquisition devices and conditions. To standardize
the dataset, Ben Graham preprocessing will be employed to normalize illumination, suppress background variation, and enhance
retinal structures by subtracting local average intensity. This will improve visibility of important pathological
features such as blood vessels, microaneurysms, exudates, and hemorrhages. Furthermore, Contrast Limited Adaptive Histogram
Equalization (CLAHE) will be applied to improve local contrast, particularly in darker or low-quality regions, making subtle
lesions more distinguishable while preventing excessive noise amplification. These preprocessing techniques will collectively
improve feature clarity, reduce inter-image variability, and help the model focus on clinically relevant patterns rather
than acquisition-related inconsistencies.

To further address class imbalance, data augmentation strategies will also be applied more aggressively to minority
classes. Transformations such as rotation, flipping, zooming, brightness adjustments, and cropping will increase sample
diversity while preserving pathological relevance. This will not only improve representation of under-sampled severity
grades but also enhance model generalization by exposing it to a broader range of retinal variations. Together,
the combination of hybrid ordinal-imbalance-aware loss functions and robust preprocessing will create a more clinically
informed and technically optimized framework for diabetic retinopathy classification, with the goal of improving
predictive accuracy across all severity levels, especially in minority and early-detection classes where accurate
diagnosis is most critical.

Building upon the baseline established using ResNet-18, which was initially trained without advanced preprocessing or 
loss function enhancements, we plan to further investigate whether alternative architectures can provide superior performance 
when combined with the proposed methodological improvements. Specifically, we will evaluate DenseNet and EfficientNet alongside 
the improved ResNet-18 framework after integrating Ben Graham preprocessing, CLAHE-based contrast enhancement, data augmentation, 
and the combined Focal-CORAL loss function. Each architecture will be trained and evaluated under the same enhanced preprocessing 
and loss function framework to ensure fair comparison. Performance will be assessed using metrics such as accuracy, F1-score, 
recall, and sensitivity across all diabetic retinopathy severity grades, with particular attention to minority and severe classes. 
In addition to predictive performance, model complexity and parameter count will also be analyzed. This comparative strategy will allow 
us to quantify not only the gains achieved through methodological improvements over the original ResNet-18 baseline, but also determine 
which architecture provides the most effective balance between classification performance and practical deployment considerations.
